\subsection{Late-time Cosmic Acceleration from Compactification --- arXiv:1811.03660}

\textbf{Authors:} J. G. Russo, P. K. Townsend

\paragraph{主な主張}
高次元の stress tensor が Strong Energy Condition と Dominant Energy Condition を満たしていても、時間依存 compactification により低次元 Einstein-frame FLRW 宇宙の永続的な power-law 加速膨張が可能であり、4次元では late-time の状態方程式が $-1/2<w<-1/3$ に制限されることを示した \cite{Russo:2018compactification}。

\paragraph{新規性}
時間依存 compactification に対する no-go 議論を精密化し、高次元 SEC は de Sitter を排除しても late-time accelerated power-law scaling までは排除しないことを具体例とともに示した点にある。

\paragraph{位置付け}
Gibbons--Maldacena--Nu\~nez 型 no-go theorem と宇宙論的 compactification の間をつなぎ、内部空間の緩やかな decompactification を伴う scaling attractor が高次元の energy condition と両立し得ることを明確にした基礎的な論文である。
